\begin{frame}{실습: 검증셋으로 k 튜닝}
  4.1절에서 ``적절한 $k$는 검증 데이터로 정한다''고 했다
  --- 직접 해보자.
  \begin{enumerate}\small
    \item 데이터를 \textbf{훈련/검증}으로 분리
    \item 여러 $k$ 각각에 대해 \textbf{검증셋 정확도} 계산
    \item 가장 높은 지점을 고른다
  \end{enumerate}
  \vskip0.6em
  \texttt{make\_moons}(scikit-learn) 500개 점, 30\% 검증 /
  70\% 훈련, $k \in \{1,3,5,7,10,15,20,30,50\}$:
  \small
  \begin{tabular}{@{}ccc@{}}
    \toprule
     $k$ & 훈련 정확도 & 검증 정확도 \\
    \midrule
    1  & \textbf{1.000} & 0.900 \\
    5  & 0.940 & 0.920 \\
    7  & 0.931 & \textbf{0.960} \\
    15 & 0.934 & 0.947 \\
    30 & 0.929 & 0.947 \\
    50 & 0.920 & 0.933 \\
    \bottomrule
  \end{tabular}
  \vskip0.8em
  $k=1$: 훈련 100\%에 가깝지만(과적합) 검증은 오히려 떨어짐 ---
  편향-분산 트레이드오프가 숫자로 드러나는 지점.
  검증 정확도는 $k$가 작을 때(분산)와 클 때(편향)
  \textbf{양쪽 끝에서 모두 낮아지는 ``U'' 모양}.
  이 표에선 $k=7$(0.960)이 최고.
\end{frame}
